% Part: lambda-calculus
% Chapter: syntax
% Section: de-bruijn

\documentclass[../../../include/open-logic-section]{subfiles}

\begin{document}

\olfileid{lam}{syn}{deb}

\olsection{مؤشر دي بروين}

التكافؤ بحسب~$\alpha$ طبيعي جدًا، لأن الحدود المتكافئة بحسب~$\alpha$
``تعني الشيء نفسه''. والواقع أنه يمكن إعطاء تركيب لحدود لامبدا لا يميز
بين الحدود التي يمكن تحويل أحدها إلى الآخر بتحويل~$\alpha$. وأشهر هذه
التراكيب يستبدل المتغيرات بـ\emph{مؤشر دي بروين} الخاص بها.

عندما نكتب~$\lambd[x][M]$، نصرّح بأن~$x$ هو معلمة الدالة، بحيث نستطيع
استعمال~$x$ في~$M$ للإحالة إلى هذه المعلمة. أما في ترميز مؤشر دي بروين
فلا أسماء للمعلمات، ويرمز إلى الإحالة إليها في متن الدالة بعدد يبين عدد
مستويات التجريد الفاصلة بين الإحالة وتجريدها. فمثلًا، تأمل
$\lambd[x][\lambd[y][y x]]$: التجريد الخارجي يربط المتغير~$x$، والتجريد
الداخلي يربط المتغير~$y$، والحد الجزئي~$y x$ يقع في مجال التجريد
الداخلي. فلا يوجد تجريد بين~$y$ وتجريده~$\lambd[y]$، ويوجد تجريد واحد
بين~$x$ وتجريده~$\lambd[x]$. لذلك نكتب~$0\, 1$ بدل~$y x$، ونكتب
$\lambd[][\lambd[][01]]$ للحد كله.

\begin{defn}
  تعرّف حدود دي بروين استقرائيًا كما يأتي:
  \begin{enumerate}
  \item $n$، حيث~$n$ أي عدد طبيعي.
  \item $PQ$، حيث~$P$ و$Q$ كلاهما حدان من حدود دي بروين.
  \item $\lambd[][N]$، حيث~$N$ حد من حدود دي بروين.
  \end{enumerate}
\end{defn}

يكون التحويل المصاغ صوريًا من حدود لامبدا العادية إلى الحدود المفهرسة
بدي بروين كما يأتي:
\begin{defn}
  \begin{align*}
    F_\Gamma(x) &= \Gamma(x) \\
    F_\Gamma(PQ) &= F_\Gamma(P)F_\Gamma(Q) \\
    F_\Gamma(\lambd[x][N]) &= \lambd[][F_{x,\Gamma}(N)]
  \end{align*}
  حيث~$\Gamma$ قائمة متغيرات مفهرسة ابتداءً من الصفر، وترمز~$\Gamma(x)$
  إلى موضع المتغير~$x$ في~$\Gamma$. فمثلًا، إذا كانت~$\Gamma$ هي
  $x,y,z$، فإن~$\Gamma(x)$ تساوي~$0$ و$\Gamma(z)$ تساوي~$2$.
  
  وترمز~$x,\Gamma$ إلى القائمة الناتجة من دفع~$x$ إلى رأس~$\Gamma$؛
  فمثلًا، متابعةً لقيمة~$\Gamma$ في المثال السابق، تكون~$w,\Gamma$ هي
  $w,x,y,z$.
\end{defn}

يستعاد حد لامبدا قياسي من حد دي بروين كما يأتي:

\begin{defn}
  \begin{align*}
    G_\Gamma(n) &= \Gamma[n] \\
    G_\Gamma(PQ) &= G_\Gamma(P) G_\Gamma(Q) \\
    G_\Gamma(\lambd[][N]) &= \lambd[x][G_{x,\Gamma}(N)]
  \end{align*}
  حيث~$\Gamma$ مرة أخرى قائمة متغيرات مفهرسة ابتداءً من الصفر، وترمز
  $\Gamma[n]$ إلى المتغير في الموضع~$n$. فمثلًا، إذا كانت~$\Gamma$ هي
  $x,y,z$، فإن~$\Gamma[1]$ هي~$y$.

  ويختار المتغير~$x$ في المعادلة الأخيرة ليكون أي متغير لا ينتمي إلى
  $\Gamma$.
\end{defn}

نعرض هنا بعض النتائج من غير برهان:

\begin{prop}
  إذا كان~$M \aconvone M'$ وكانت~$\Gamma$ أي قائمة تحتوي~$\FV{M}$، فإن
  $F_\Gamma(M) \eqs F_\Gamma(M')$.
\end{prop}

\end{document}
